% =============================================================================
%  1-Introduccion.tex — Contexto territorial y planteamiento del problema
% =============================================================================

\section{Introducción y contexto territorial}

La Zona Metropolitana del Valle de México (ZMVM) constituye uno de los conglomerados
urbanos más complejos y densamente poblados del continente.  Históricamente, la
planeación y expansión de su red de transporte masivo y semimasivo ha seguido un
patrón predominantemente radial, orientado a conectar las periferias con el núcleo
central de la actividad económica en la Ciudad de México
\citep{padilla2022expansion, molinero2002}.  Esta morfología centro-periferia genera
una alta vulnerabilidad funcional en el sistema completo: obliga a recorridos
ineficientes, incrementa los tiempos de viaje de la población trabajadora y satura
las líneas troncales en puntos de transbordo críticos
\citep{semovi2020diagnostico}.

El estudio cuantitativo y espacial del transporte público a escala de megalópolis
enfrenta un obstáculo recurrente: la fragmentación institucional de la información
geoespacial.  Si bien existen iniciativas de datos abiertos gubernamentales y
especificaciones estandarizadas como los \textit{General Transit Feed Specification}
(GTFS) proveídos por la Secretaría de Movilidad (SEMOVI), estos archivos presentan
severas discontinuidades cuando la red cruza las fronteras políticas entre la Ciudad
de México y los municipios conurbados del Estado de México.  Los portales oficiales
suelen ofrecer geometrías desactualizadas, tramos rotos, atributos inconsistentes
o, en el peor de los casos, la exclusión total de sistemas periféricos críticos como
el Mexibús o el Trolebús elevado \citep{hadas2013assessing}.

Esta fricción de entrada contradice el principio de acceso a la información: la
apertura formal de un dato gubernamental no equivale a su democratización si este
requiere habilidades avanzadas de ingeniería de software para ser interpretado y
explotado científicamente.  Se vuelve indispensable, por tanto, contar con una
infraestructura intermedia que unifique, limpie y exponga la red de transporte
como un recurso directamente consumible por herramientas de análisis espacial.
